\documentclass[11pt]{article}
% Preámbulo compartido por los artefactos Scrum del proyecto
% Proyecto Agile — Spanish Braille Application (EPN)
\usepackage[utf8]{inputenc}
\usepackage[T1]{fontenc}
\usepackage[spanish]{babel}
\usepackage[a4paper,margin=2.3cm]{geometry}
\usepackage{amsmath}
\usepackage{amssymb}
\usepackage{xcolor}
\usepackage{booktabs}
\usepackage{longtable}
\usepackage{array}
\usepackage{colortbl}
\usepackage{enumitem}
\usepackage{fancyhdr}
\usepackage{titlesec}
\usepackage{tcolorbox}
\usepackage{pgfplots}
\pgfplotsset{compat=1.18}
\usepackage{hyperref}

\definecolor{epnblue}{HTML}{2E4C9C}
\definecolor{epnbluelight}{HTML}{4B6FCB}
\definecolor{epnrow}{HTML}{EDF1FB}
\definecolor{epngray}{HTML}{555555}

\hypersetup{
  colorlinks=true,
  linkcolor=epnblue,
  urlcolor=epnblue,
  citecolor=epnblue,
  pdfcreator={XeLaTeX},
}

\titleformat{\section}
  {\normalfont\Large\bfseries\color{epnblue}}{\thesection}{0.6em}{}
\titleformat{\subsection}
  {\normalfont\large\bfseries\color{epnbluelight}}{\thesubsection}{0.6em}{}

\pagestyle{fancy}
\fancyhf{}
\renewcommand{\headrulewidth}{0.4pt}
\renewcommand{\footrulewidth}{0.2pt}
\fancyfoot[C]{\thepage}

\newcommand{\artefactoTitulo}[2]{%
  \begin{titlepage}
    \centering
    \vspace*{1.2cm}
    {\large Escuela Politécnica Nacional\par}
    {\normalsize Facultad de Ingeniería de Sistemas\par}
    {\normalsize Construcción y Evolución de Software --- Primer Bimestre\par}
    \vspace{1.4cm}
    \begin{tcolorbox}[colback=epnblue,colframe=epnblue,coltext=white,arc=2mm,boxrule=0pt,width=0.92\textwidth,halign=center]
      {\Huge\bfseries #1\par}
    \end{tcolorbox}
    \vspace{0.5cm}
    {\Large #2\par}
    \vspace{1.6cm}
    {\large\bfseries Proyecto Agile --- Spanish Braille Application\par}
    \vspace{0.3cm}
    {\normalsize Traducción bidireccional español \(\leftrightarrow\) Braille Unicode\par}
    \vfill
    \begin{tabular}{ll}
      \textbf{Product Owner:}  & José Castro \\
      \textbf{Scrum Master:}   & Victor Aveiga \\
      \textbf{Equipo de Desarrollo:} & Javier Angulo, Elian Caizapanta, Erick Costa, Emily Aumala \\
    \end{tabular}
    \vspace{0.8cm}
    \par\noindent\rule{0.5\textwidth}{0.4pt}\par
    \vspace{0.3cm}
    {\normalsize Documento preparado por: Erick Costa (Equipo de Desarrollo)\par}
    {\normalsize \today\par}
  \end{titlepage}
  \fancyhead[L]{Proyecto Agile --- Spanish Braille App}
  \fancyhead[R]{#1}
}


\title{Product Backlog}
\author{}
\date{}

\begin{document}
\artefactoTitulo{Product Backlog}{Lista priorizada de requisitos del producto}

\section{Introducción}
Este documento contiene el \textbf{Product Backlog} del proyecto \textit{Spanish Braille Application},
una aplicación web (Spring Boot + Thymeleaf) que traduce texto en español a Braille Unicode y viceversa,
incluyendo acentos, la letra \textbf{ñ}, números, mayúsculas, signos de puntuación y una variante de
\textbf{Braille espejo} para impresión en relieve.

El Backlog es mantenido por el \textbf{Product Owner (José Castro)} y refleja las Historias de Usuario (HU)
priorizadas para los dos sprints ejecutados hasta la fecha, derivadas del documento de lógica de negocio
(\texttt{docs/design/roadmap-alfabeto-braille.md}) y verificadas contra la implementación real
(\texttt{src/main/java/.../service}, \texttt{controller}) y las pruebas unitarias del repositorio.

\section{Ítems del Product Backlog}

\renewcommand{\arraystretch}{1.35}
\begin{longtable}{|>{\raggedright\arraybackslash}p{1.6cm}|p{6.7cm}|>{\centering\arraybackslash}p{1.2cm}|>{\centering\arraybackslash}p{1.6cm}|>{\centering\arraybackslash}p{1.1cm}|p{2.5cm}|}
\hline
\rowcolor{epnblue}
\textcolor{white}{\textbf{ID}} & \textcolor{white}{\textbf{Historia de Usuario}} & \textcolor{white}{\textbf{Pts}} & \textcolor{white}{\textbf{Prioridad}} & \textcolor{white}{\textbf{Sprint}} & \textcolor{white}{\textbf{Estado}} \\
\hline
\endfirsthead
\hline
\rowcolor{epnblue}
\textcolor{white}{\textbf{ID}} & \textcolor{white}{\textbf{Historia de Usuario}} & \textcolor{white}{\textbf{Pts}} & \textcolor{white}{\textbf{Prioridad}} & \textcolor{white}{\textbf{Sprint}} & \textcolor{white}{\textbf{Estado}} \\
\hline
\endhead

\rowcolor{epnrow}
01H1 & Como usuario, quiero transcribir el abecedario español (a--z, series 1.\textsuperscript{a}, 2.\textsuperscript{a} y 3.\textsuperscript{a}) a Braille Unicode, para poder generar texto Braille legible a partir de texto plano. & 5 & Alta & 1 & \cellcolor{green!25}Hecho \\
\hline
02H2 & Como usuario, quiero convertir vocales acentuadas (á, é, í, ó, ú) y la letra ñ a su representación Braille específica, para transcribir correctamente palabras en español con tildes y diéresis. & 3 & Media & 1 & \cellcolor{green!25}Hecho \\
\hline
\rowcolor{epnrow}
03H3 & Como usuario, quiero transcribir números (0--9) anteponiendo el signo numérico Braille (puntos 3-4-5-6), para representar cifras de forma correcta según el estándar Braille español. & 2 & Alta & 1 & \cellcolor{green!25}Hecho \\
\hline
04H4 & Como usuario, quiero que las letras mayúsculas se antepongan con el indicador Braille de mayúscula (puntos 4-6), para distinguir mayúsculas de minúsculas en el texto transcrito. & 2 & Media & 1 & \cellcolor{green!25}Hecho \\
\hline
\rowcolor{epnrow}
05H5 & Como usuario, quiero convertir signos de puntuación básicos (\texttt{. , ; : ¿ ? ¡ !}) a su representación Braille, para transcribir oraciones completas con su puntuación original. & 3 & Media & 2 & \cellcolor{green!25}Hecho \\
\hline
06H6 & Como usuario con discapacidad visual o docente de braille, quiero generar señalética Braille imprimible a partir de un texto, para producir material accesible en relieve. & 7 & Alta & 2 & \cellcolor{yellow!40}Parcial \\
\hline
\rowcolor{epnrow}
07H7 & Como usuario, quiero transcribir un patrón Braille Unicode de vuelta a español (conversión inversa), para poder leer e interpretar contenido escrito originalmente en Braille. & 5 & Alta & 2 & \cellcolor{green!25}Hecho \\
\hline
\end{longtable}

\vspace{0.3cm}
\noindent\textbf{Story points totales:} 27 pts

\noindent\textbf{Leyenda de estado:} \colorbox{green!25}{Hecho} = funcionalidad implementada y cubierta por pruebas unitarias
verificables en \texttt{src/test/java}; \colorbox{yellow!40}{Parcial} = el equipo (Elian Caizapanta) reporta
la exportación de señalética como completada y la vista permite imprimir el resultado en espejo
(botón ``Imprimir espejo'' en \texttt{result-braille.html}), pero la rama de este repositorio aún no
refleja una exportación dedicada a PNG/PDF independiente del navegador --- pendiente de confirmar
sincronización antes de la entrega (ver Sprint Backlog).

\section{Criterios de priorización}
El backlog se priorizó combinando:
\begin{itemize}[leftmargin=1.4cm]
  \item \textbf{Valor para el usuario final} (accesibilidad: transcripción bidireccional funcional primero).
  \item \textbf{Dependencias técnicas}: el abecedario base (01H1) es prerequisito de acentos, mayúsculas,
        números y puntuación, por lo que se ubicó primero en el Sprint 1.
  \item \textbf{Complejidad/Riesgo}: 06H6 (señalética exportable) es la historia de mayor esfuerzo (7 pts)
        y mayor incertidumbre técnica, por lo que se dejó al final del Sprint 2.
\end{itemize}

\section{Refinamiento continuo}
Este Product Backlog es un artefacto vivo: el Product Owner (José Castro), junto con el equipo de desarrollo,
lo revisa al inicio de cada Sprint Planning para reordenar prioridades, dividir historias grandes
(p. ej. 06H6 podría dividirse en ``renderizar señalética'' + ``exportar a PNG'' + ``exportar a PDF''
en una futura iteración) y agregar nuevas historias que surjan de la retroalimentación de usuarios.

\end{document}
